\documentclass[11pt,a4paper]{article}
\usepackage[utf8]{inputenc}
\usepackage{amsmath,amssymb,amsfonts}
\usepackage{graphicx}
\usepackage{geometry}
\usepackage{enumitem}
\usepackage{fancyhdr}

\geometry{top=1in, bottom=1in, left=0.8in, right=0.8in}
\pagestyle{fancy}
\fancyhf{}
\rhead{Max Marks: 50}
\lhead{Exam: World War II & Decolonization}
\rfoot{Page \thepage}

\begin{document}

\begin{center}
    {\LARGE \textbf{ World War II & Decolonization }} \\[0.5em]
    \textbf{Total Maximum Marks:} 50 \quad | \quad \textbf{Total Questions:} 2
\end{center}

\vspace{1em}
\hrule
\vspace{1.5em}

\noindent \textbf{Instructions:} Answer all questions carefully. Show all mathematical derivations, intermediate steps, and calculations where applicable.

\vspace{1.5em}

\begin{enumerate}[label=\textbf{Q\arabic*.}, leftmargin=2em]

    \item \textbf{[25 Marks]} Explain how the surrender of the Axis powers in 1945 led to the creation of the United Nations. In your answer, discuss the key provisions of the UN Charter that were designed to prevent future global conflicts and evaluate the extent to which the UN has succeeded in maintaining international peace since its founding.
    
    \vspace{1em}

    \item \textbf{[25 Marks]} Analyze the Indian Independence Act of 1947 as a turning point in the process of decolonization. Compare its impact on British imperial policy with that of another decolonization case study of your choice (e.g., Ghana 1957, Algeria 1962, or Indonesia 1949). Discuss the political, social, and economic consequences of partition for the newly formed states of India and Pakistan.
    
    \vspace{1em}

\end{enumerate}

\end{document}